\documentclass[11pt]{article}
\usepackage[utf8]{inputenc}
\usepackage{graphicx}
\usepackage{amsmath}
\usepackage{amssymb}
\usepackage{mathrsfs}
\usepackage{pgfornament}
\usepackage[many]{tcolorbox}
\usepackage[margin = 0.5in]{geometry}
\usepackage{collectbox}
\usepackage{mdframed}
\usepackage{xcolor}

\newcommand\textbox[1]{%
  \parbox{.333\textwidth}{#1}%
}

\linespread{1.5}

\makeatletter
\newcommand{\mybox}{%
    \collectbox{%
        \setlength{\fboxsep}{2pt}%
        \fbox{\BOXCONTENT}%
    }%
}
\makeatother

\usepackage{listings}
\usepackage{color}

\definecolor{dkgreen}{rgb}{0,0.6,0}
\definecolor{gray}{rgb}{0.5,0.5,0.5}
\definecolor{mauve}{rgb}{0.58,0,0.82}

\lstset{frame=tb,
  language=R,
  aboveskip=3mm,
  belowskip=3mm,
  showstringspaces=false,
  columns=flexible,
  basicstyle={\small\ttfamily},
  numbers=none,
  numberstyle=\tiny\color{gray},
  keywordstyle=\color{blue},
  commentstyle=\color{dkgreen},
  stringstyle=\color{mauve},
  breaklines=true,
  breakatwhitespace=true,
  tabsize=3
}
\title{MAT 586: Lab 8}
\begin{document}

%%
%% Title 
%%
\begin{center}
\vspace{-0.6cm}
\hrulefill \\
\vspace{-0.7cm}
\hrulefill

\vspace{-0.5cm}
\begin{flushleft}
\noindent\textbox{\pgfornament[width = 9mm, color = black]{39}\hfill}\textbox{\hfil \textbf{\begin{large}MAT 586 \end{large}}\hfil}\textbox{\hfill \pgfornament[width = 9mm, color = black]{40}}
\end{flushleft}
\vspace{-5mm}
\begin{large}
\textbf{Lab 8} \\
\end{large}

\vspace{5mm}

\begin{flushleft}
\noindent\textbox{\pgfornament[width = 9mm, color = black, symmetry = h]{39}\hfill}\textbox{\hfil \textbf{\begin{large}Nolan R. H. Gagnon\end{large}}\hfill}\textbox{\hfill \pgfornament[width = 9mm, color = black, symmetry = h]{40}}
\end{flushleft}

\vspace{-0.65cm}
\hrulefill \\
\vspace{-0.7cm}
\hrulefill
\end{center}

\begin{center}
\pgfornament[width = 75mm,
             color = black]{89}
\end{center}

\vspace{5mm}

\noindent\mybox{\colorbox{lightgray}{\textbf{\# 1}}}\hspace{3mm}\hrulefill 

\noindent The continuous-time logistic differential equation is
$$
\frac{dy}{dt} = ry\left(1 - \frac{y}{K}\right).
$$
If the initial condition is $y(0) = y_0$, then the solution to the differential equation is 
$$
y(t) = \frac{Ky_0e^{rt}}{K + y_0\left(e^{rt}-1\right)}.
$$

\begin{figure}[h!]
\centering
  \includegraphics[width = 0.6\textwidth]{ls_sol.pdf}	
  \caption{Estimates of solution to logistic differential equation}
  \label{fig:logistic_est}
\end{figure}

\noindent Figure 1 shows ``average'' and least-squares solutions to the differential equation with unknown $K$ and $r$ randomly chosen from unknown probability distributions. The least-squares estimates of $K$ and $r$ were $\hat{K} = 16.9370488$ and $\hat{r} = 0.6654223$. These estimates (and the figure above) were generated with the following R code.

\pagebreak

\noindent\mybox{\colorbox{yellow}{\textbf{lab8\_solution.R}}}\hspace{3mm}\hrulefill
\begin{lstlisting}
source("log_sol.R") ## Contains function with params to estimate with least-squares optimization.
source("lab8data.R") ## Average data  
library('Hmisc') ## Gives access to errbar function

par.v <- optim(c(1, 1), log_sol, t.v = t.v, y.v = avgY.v)$par

## Since y(0) = y0, it must be the case
## that y0 = avgY.v[1] = (20*y0) / 20
y0 <- avgY.v[1]

## Least-squares estimate of K
K_hat <- par.v[1]

## Least-squares estimate of r
r_hat <- par.v[2]

## Calculate the least-squares solution for y(t)
y.v <- (K_hat * y0 * exp(r_hat * t.v)) / (K_hat + y0 * (exp(r_hat * t.v) - 1))

## Plot the average solution to the logistic diff. eq.
## along with the solution found by least-squares 
## optimization.
errbar(t.v, avgY.v, avgY.v + sdY.v, avgY.v - sdY.v)
lines(t.v, y.v, col = "red", lwd = 3)
title(main = "Estimates of solution to logistic diff. eq. with params. K and r")
legend("bottomright", legend = "Least-squares Sol.", lty = "solid", col = "red")
\end{lstlisting}

\vspace{10mm}

\noindent\mybox{\colorbox{yellow}{\textbf{log\_sol.R}}}\hspace{3mm}\hrulefill
\begin{lstlisting}
log_sol <- function(beta.v, t.v, y.v)
{
	## Get parameters from beta vector
	K <- beta.v[1]
	r <- beta.v[2]
	
	## Get y0 from average y vector
	y0 <- y.v[1]	

	## Compute predicted values	
	yhat.v <- (K * y0 * exp(r * t.v)) / (K + y0 * (exp(r * t.v) - 1))
	
	## Compute residuals
	res.v <- y.v - yhat.v

	## Return sum of squared residuals	
	return(sum(res.v ^ 2))
}
\end{lstlisting}

\hfill$\blacksquare$


\vfill
\begin{center}
\pgfornament[width = 75mm,
             color = black]{89}
\end{center}
\begin{center}
\vspace{-0.6cm}
\hrulefill \\
\vspace{-0.7cm}
\hrulefill
\end{center}

\end{document}
